% GENERATED FILE. Edit canonical lessons, then run npm run generate:books.
% canonical-source-hash: fnv1a64:056024e92f76081f
% canonical-lessons: JA-C331-shakai, JA-C331-shimin, JA-C331-juumin, JA-C331-seiji, JA-C331-houritsu

\chapter{Society}
\label{ch:ja-society}
\hlchaptermodality{331}

\begin{chapteropening}
\textbf{By the end of this chapter:} \emph{I can give my opinion about how society is changing.}

Complete the last lesson of chapter 331: I can give my opinion about how society is changing.
\end{chapteropening}

\section[\texorpdfstring{\ja{しゃかい} (shakai)}{しゃかい (shakai)}]{\ja{しゃかい} (shakai) — society}

\label{lesson:JA-C331-shakai}



\begin{quote}
Before the new one: say the Japanese for to line up neatly, then the Japanese for to revise.
\end{quote}

\subsection*{You'll want to know: \ja{しゃかい}}
\textbf{\ja{しゃかい}} — \emph{shakai} — \textquotedblleft{}society\textquotedblright{}.

\begin{quote}
\textbf{\ja{しゃかいは} \ja{かわったと} \ja{おもいます。}} — \emph{shakai wa kawatta to omoimasu} — \textquotedblleft{}I think society has changed.\textquotedblright{}
\end{quote}

\begin{cousinweb}
\textbf{\ja{しゃかい}} (\emph{shakai}) is Sino-Japanese (kango): two Chinese characters read with their on'yomi, readings borrowed from Chinese. They are \emph{sha}, \textquotedblleft{}association\textquotedblright{} (Middle Chinese \emph{dzyaeX}) and \emph{kai}, \textquotedblleft{}meeting\textquotedblright{} (\emph{hwajH}) (KANJIDIC2; Middle Chinese in Baxter's transcription).
\end{cousinweb}

\begin{grammarlens}[title={What you've built}]
One more word for giving reasons and opinions.
\end{grammarlens}

\subsection*{Your turn}
\begin{itemize}
\raggedright
  \item \emph{Say it:} \emph{shakai}
  \item \emph{Say it:} \emph{shakai}, once more
  \item \emph{Say it:} \emph{shakai wa kawatta to omoimasu} — I think society has changed.
  \item \emph{Recall:} say the Japanese for to line up neatly, then the Japanese for to revise, then say \emph{shakai} again
\end{itemize}

\begin{tcolorbox}[breakable,colback=teal!4,colframe=teal!35!black,title={Before you move on}]
What does \ja{しゃかい} mean? (\textquotedblleft{}Society\textquotedblright{}.) Say it once more.
\end{tcolorbox}
\section[\texorpdfstring{\ja{しみん} (shimin)}{しみん (shimin)}]{\ja{しみん} (shimin) — a citizen}

\label{lesson:JA-C331-shimin}



\begin{quote}
Before the new one: say the Japanese for to revise, then the Japanese for society.
\end{quote}

\subsection*{You'll want to know: \ja{しみん}}
\textbf{\ja{しみん}} — \emph{shimin} — \textquotedblleft{}a citizen\textquotedblright{}.

\begin{quote}
\textbf{\ja{しみんの} \ja{いけんを} \ja{きいた} \ja{ほうが} \ja{いいです。}} — \emph{shimin no iken o kiita hō ga ii desu} — \textquotedblleft{}They ought to listen to what the citizens think.\textquotedblright{}
\end{quote}

\begin{cousinweb}
\textbf{\ja{しみん}} (\emph{shimin}) is Sino-Japanese (kango): two Chinese characters read with their on'yomi, readings borrowed from Chinese. They are \emph{shi}, \textquotedblleft{}city\textquotedblright{} (Middle Chinese \emph{dzyiX}) and \emph{min}, \textquotedblleft{}people\textquotedblright{} (\emph{mjin}) (KANJIDIC2; Middle Chinese in Baxter's transcription).
\end{cousinweb}

\begin{grammarlens}[title={What you've built}]
One more word for giving reasons and opinions.
\end{grammarlens}

\subsection*{Your turn}
\begin{itemize}
\raggedright
  \item \emph{Say it:} \emph{shimin}
  \item \emph{Say it:} \emph{shimin}, once more
  \item \emph{Say it:} \emph{shimin no iken o kiita hō ga ii desu} — They ought to listen to what the citizens think.
  \item \emph{Recall:} say the Japanese for to revise, then the Japanese for society, then say \emph{shimin} again
\end{itemize}

\begin{tcolorbox}[breakable,colback=teal!4,colframe=teal!35!black,title={Before you move on}]
What does \ja{しみん} mean? (\textquotedblleft{}A citizen\textquotedblright{}.) Say it once more.
\end{tcolorbox}
\section[\texorpdfstring{\ja{じゅうみん} (jūmin)}{じゅうみん (jūmin)}]{\ja{じゅうみん} (jūmin) — a resident; the local people}

\label{lesson:JA-C331-juumin}



\begin{quote}
Before the new one: say the Japanese for society, then the Japanese for a citizen.
\end{quote}

\subsection*{You'll want to know: \ja{じゅうみん}}
\textbf{\ja{じゅうみん}} — \emph{jūmin} — \textquotedblleft{}a resident; the local people\textquotedblright{}.

\begin{quote}
\textbf{\ja{じゅうみんが} \ja{はんたいしたので、ちゅうしに} \ja{なりました。}} — \emph{jūmin ga hantai shita node, chūshi ni narimashita} — \textquotedblleft{}The residents were against it, so it was called off.\textquotedblright{}
\end{quote}

\begin{cousinweb}
\textbf{\ja{じゅうみん}} (\emph{jūmin}) is Sino-Japanese (kango): two Chinese characters read with their on'yomi, readings borrowed from Chinese. They are \emph{jū}, \textquotedblleft{}dwell\textquotedblright{} (Middle Chinese \emph{drjuH}) and \emph{min}, \textquotedblleft{}people\textquotedblright{} (\emph{mjin}) (KANJIDIC2; Middle Chinese in Baxter's transcription).
\end{cousinweb}

\begin{grammarlens}[title={What you've built}]
One more word for giving reasons and opinions.
\end{grammarlens}

\subsection*{Your turn}
\begin{itemize}
\raggedright
  \item \emph{Say it:} \emph{jūmin}
  \item \emph{Say it:} \emph{jūmin}, once more
  \item \emph{Say it:} \emph{jūmin ga hantai shita node, chūshi ni narimashita} — The residents were against it, so it was called off.
  \item \emph{Recall:} say the Japanese for society, then the Japanese for a citizen, then say \emph{jūmin} again
\end{itemize}

\begin{tcolorbox}[breakable,colback=teal!4,colframe=teal!35!black,title={Before you move on}]
What does \ja{じゅうみん} mean? (\textquotedblleft{}A resident; the local people\textquotedblright{}.) Say it once more.
\end{tcolorbox}
\section[\texorpdfstring{\ja{せいじ} (seiji)}{せいじ (seiji)}]{\ja{せいじ} (seiji) — politics}

\label{lesson:JA-C331-seiji}



\begin{quote}
Before the new one: say the Japanese for a citizen, then the Japanese for a resident.
\end{quote}

\subsection*{You'll want to know: \ja{せいじ}}
\textbf{\ja{せいじ}} — \emph{seiji} — \textquotedblleft{}politics\textquotedblright{}.

\begin{quote}
\textbf{\ja{しょうらいの} \ja{ために、せいじを} \ja{べんきょうしたいです。}} — \emph{shōrai no tame ni, seiji o benkyō shitai desu} — \textquotedblleft{}I want to study politics for my future.\textquotedblright{}
\end{quote}

\begin{cousinweb}
\textbf{\ja{せいじ}} (\emph{seiji}) is Sino-Japanese (kango): two Chinese characters read with their on'yomi, readings borrowed from Chinese. They are \emph{sei}, \textquotedblleft{}politics\textquotedblright{} (Middle Chinese \emph{tsyengH}) and \emph{ji}, \textquotedblleft{}reign\textquotedblright{} (\emph{dri}) (KANJIDIC2; Middle Chinese in Baxter's transcription).
\end{cousinweb}

\begin{grammarlens}[title={What you've built}]
One more word for giving reasons and opinions.
\end{grammarlens}

\subsection*{Your turn}
\begin{itemize}
\raggedright
  \item \emph{Say it:} \emph{seiji}
  \item \emph{Say it:} \emph{seiji}, once more
  \item \emph{Say it:} \emph{shōrai no tame ni, seiji o benkyō shitai desu} — I want to study politics for my future.
  \item \emph{Recall:} say the Japanese for a citizen, then the Japanese for a resident, then say \emph{seiji} again
\end{itemize}

\begin{tcolorbox}[breakable,colback=teal!4,colframe=teal!35!black,title={Before you move on}]
What does \ja{せいじ} mean? (\textquotedblleft{}Politics\textquotedblright{}.) Say it once more.
\end{tcolorbox}
\section[\texorpdfstring{\ja{ほうりつ} (hōritsu)}{ほうりつ (hōritsu)}]{\ja{ほうりつ} (hōritsu) — a law}

\label{lesson:JA-C331-houritsu}



\begin{quote}
Before the new one: say the Japanese for a resident, then the Japanese for politics.
\end{quote}

\subsection*{You'll want to know: \ja{ほうりつ}}
\textbf{\ja{ほうりつ}} — \emph{hōritsu} — \textquotedblleft{}a law\textquotedblright{}.

\begin{quote}
\textbf{\ja{ほうりつが} \ja{かわったので、ぜいきんが} \ja{たかく} \ja{なりました。}} — \emph{hōritsu ga kawatta node, zeikin ga takaku narimashita} — \textquotedblleft{}The law changed, so taxes went up.\textquotedblright{}
\end{quote}

\begin{cousinweb}
\textbf{\ja{ほうりつ}} (\emph{hōritsu}) is Sino-Japanese (kango): two Chinese characters read with their on'yomi, readings borrowed from Chinese. They are \emph{hō}, \textquotedblleft{}law\textquotedblright{} (Middle Chinese \emph{pjop}) and \emph{ritsu}, \textquotedblleft{}regulation\textquotedblright{} (\emph{lwit}) (KANJIDIC2; Middle Chinese in Baxter's transcription).
\end{cousinweb}

\begin{grammarlens}[title={What you've built}]
One more word for giving reasons and opinions.
\end{grammarlens}

\subsection*{Your turn}
\begin{itemize}
\raggedright
  \item \emph{Say it:} \emph{hōritsu}
  \item \emph{Say it:} \emph{hōritsu}, once more
  \item \emph{Say it:} \emph{hōritsu ga kawatta node, zeikin ga takaku narimashita} — The law changed, so taxes went up.
  \item \emph{Recall:} say the Japanese for a resident, then the Japanese for politics, then say \emph{hōritsu} again
\end{itemize}

\begin{tcolorbox}[breakable,colback=teal!4,colframe=teal!35!black,title={Before you move on}]
What does \ja{ほうりつ} mean? (\textquotedblleft{}A law\textquotedblright{}.) Say it once more.
\end{tcolorbox}
